% GENERATED FILE. Edit canonical lessons, then run npm run generate:books.
% canonical-source-hash: fnv1a64:116b2e72f01f1a47
% canonical-lessons: FA-C482-delkhor, FA-C482-ranjidan, FA-C482-azorde, FA-C482-kine, FA-C482-sharmande

\chapter{Hurt Feelings}
\label{ch:fa-hurt-feelings}
\hlchaptermodality{482}

\begin{chapteropening}
\textbf{By the end of this chapter:} \emph{I can say when someone is hurt or offended.}

Complete the last lesson of chapter 482: I can say when someone is hurt or offended.
\end{chapteropening}

\section[\texorpdfstring{\fa{دلخور} (delkhor)}{دلخور (delkhor)}]{\fa{دلخور} (delkhor) — upset (with someone)}

\label{lesson:FA-C482-delkhor}



\begin{quote}
Before the new one: say the Persian for calm, then the Persian for to get angry.
\end{quote}

\subsection*{You'll want to know: \fa{دلخور}}
\textbf{\fa{دلخور}} (\emph{delkhor}) — \textquotedblleft{}upset (with someone)\textquotedblright{}.

\begin{quote}
\fa{خواهرم} \fa{از} \fa{حرف} \fa{من} \fa{دلخور} \fa{است}.
\end{quote}

\begin{quote}
\emph{khâharam az harf-e man delkhor ast.}
\end{quote}

\begin{quote}
My sister is upset at what I said.
\end{quote}

The cause takes \emph{az}: \emph{az kasi delkhor budan}. Built on \emph{del} 'heart' and \emph{khor}, the stem of \emph{khordan}.

\begin{cousinweb}
\emph{delkhor} is built on two native Persian pieces, usually explained as \emph{del}, 'heart', and \emph{khor}, the present stem of \emph{khordan}, 'to eat': a heart that is eaten at. EtymDB traces \emph{del} to Middle Persian (EtymDB).
\end{cousinweb}

\begin{grammarlens}[title={What you've built}]
One more word for telling, explaining and describing at B1.
\end{grammarlens}

\subsection*{Your turn}
\begin{itemize}
\raggedright
  \item \emph{Say it:} \emph{delkhor}
  \item \emph{Say it:} \emph{delkhor}, once more
  \item \emph{Say it:} a sentence of your own with \emph{delkhor}
  \item \emph{Recall:} say the Persian for calm, then the Persian for to get angry, then say \emph{delkhor} again
\end{itemize}

\begin{tcolorbox}[breakable,colback=teal!4,colframe=teal!35!black,title={Before you move on}]
What does \fa{دلخور} mean? (\textquotedblleft{}Upset (with someone)\textquotedblright{}.) Say it once more.
\end{tcolorbox}
\section[\texorpdfstring{\fa{رنجیدن} (ranjidan)}{رنجیدن (ranjidan)}]{\fa{رنجیدن} (ranjidan) — to take offence; to be hurt}

\label{lesson:FA-C482-ranjidan}



\begin{quote}
Before the new one: say the Persian for to get angry, then the Persian for upset.
\end{quote}

\subsection*{You'll want to know: \fa{رنجیدن}}
\textbf{\fa{رنجیدن}} (\emph{ranjidan}) — \textquotedblleft{}to take offence; to be hurt\textquotedblright{}.

\begin{quote}
\fa{او} \fa{از} \fa{شوخی} \fa{دوستش} \fa{رنجید}.
\end{quote}

\begin{quote}
\emph{u az shukhi-ye dustash ranjid.}
\end{quote}

\begin{quote}
He was hurt by his friend's joke.
\end{quote}

Present stem \emph{ranj-}: \emph{miranjad} 'takes offence'. The noun \emph{ranj} means 'suffering, toil'.

\begin{cousinweb}
\emph{ranjidan} is built on the native Persian noun \emph{ranj}, 'suffering, toil', with the verb ending \emph{-idan}. EtymDB has no entry for it, and the older history of \emph{ranj} is not settled.
\end{cousinweb}

\begin{grammarlens}[title={What you've built}]
One more word for telling, explaining and describing at B1.
\end{grammarlens}

\subsection*{Your turn}
\begin{itemize}
\raggedright
  \item \emph{Say it:} \emph{ranjidan}
  \item \emph{Say it:} \emph{ranjidan}, once more
  \item \emph{Say it:} a sentence of your own with \emph{ranjidan}
  \item \emph{Recall:} say the Persian for to get angry, then the Persian for upset, then say \emph{ranjidan} again
\end{itemize}

\begin{tcolorbox}[breakable,colback=teal!4,colframe=teal!35!black,title={Before you move on}]
What does \fa{رنجیدن} mean? (\textquotedblleft{}To take offence; to be hurt\textquotedblright{}.) Say it once more.
\end{tcolorbox}
\section[\texorpdfstring{\fa{آزرده} (âzorde)}{آزرده (âzorde)}]{\fa{آزرده} (âzorde) — hurt (in one's feelings), offended}

\label{lesson:FA-C482-azorde}



\begin{quote}
Before the new one: say the Persian for upset, then the Persian for to take offence.
\end{quote}

\subsection*{You'll want to know: \fa{آزرده}}
\textbf{\fa{آزرده}} (\emph{âzorde}) — \textquotedblleft{}hurt (in one's feelings), offended\textquotedblright{}.

\begin{quote}
\fa{او} \fa{از} \fa{رفتار} \fa{سرد} \fa{همکارش} \fa{آزرده} \fa{شد}.
\end{quote}

\begin{quote}
\emph{u az raftâr-e sard-e hamkârash âzorde shod.}
\end{quote}

\begin{quote}
She was hurt by her colleague's cold manner.
\end{quote}

The past participle of \emph{âzordan} 'to hurt, to offend' (present stem \emph{âzâr-}). \emph{âzâr} also means 'harm'.

\begin{cousinweb}
\emph{âzorde} is built on the native Persian verb \emph{âzordan}, 'to hurt, to offend': it is that verb's past participle. EtymDB has no entry for it, and its older form is not settled.
\end{cousinweb}

\begin{grammarlens}[title={What you've built}]
One more word for telling, explaining and describing at B1.
\end{grammarlens}

\subsection*{Your turn}
\begin{itemize}
\raggedright
  \item \emph{Say it:} \emph{âzorde}
  \item \emph{Say it:} \emph{âzorde}, once more
  \item \emph{Say it:} a sentence of your own with \emph{âzorde}
  \item \emph{Recall:} say the Persian for upset, then the Persian for to take offence, then say \emph{âzorde} again
\end{itemize}

\begin{tcolorbox}[breakable,colback=teal!4,colframe=teal!35!black,title={Before you move on}]
What does \fa{آزرده} mean? (\textquotedblleft{}Hurt (in one's feelings), offended\textquotedblright{}.) Say it once more.
\end{tcolorbox}
\section[\texorpdfstring{\fa{کینه} (kine)}{کینه (kine)}]{\fa{کینه} (kine) — a grudge; spite}

\label{lesson:FA-C482-kine}



\begin{quote}
Before the new one: say the Persian for to take offence, then the Persian for hurt.
\end{quote}

\subsection*{You'll want to know: \fa{کینه}}
\textbf{\fa{کینه}} (\emph{kine}) — \textquotedblleft{}a grudge; spite\textquotedblright{}.

\begin{quote}
\fa{او} \fa{از} \fa{کسی} \fa{کینه} \fa{به} \fa{دل} \fa{نمیگیرد}.
\end{quote}

\begin{quote}
\emph{u az kasi kine be del nemigirad.}
\end{quote}

\begin{quote}
He doesn't hold a grudge against anyone.
\end{quote}

\emph{kine dâshtan} or \emph{kine be del gereftan} — to bear a grudge.

\begin{cousinweb}
\emph{kine} is a native Persian word, usually traced to Middle Persian \emph{kēn}, 'hatred, revenge'. EtymDB has no entry for it, so the link is hedged.
\end{cousinweb}

\begin{grammarlens}[title={What you've built}]
One more word for telling, explaining and describing at B1.
\end{grammarlens}

\subsection*{Your turn}
\begin{itemize}
\raggedright
  \item \emph{Say it:} \emph{kine}
  \item \emph{Say it:} \emph{kine}, once more
  \item \emph{Say it:} a sentence of your own with \emph{kine}
  \item \emph{Recall:} say the Persian for to take offence, then the Persian for hurt, then say \emph{kine} again
\end{itemize}

\begin{tcolorbox}[breakable,colback=teal!4,colframe=teal!35!black,title={Before you move on}]
What does \fa{کینه} mean? (\textquotedblleft{}A grudge; spite\textquotedblright{}.) Say it once more.
\end{tcolorbox}
\section[\texorpdfstring{\fa{شرمنده} (sharmande)}{شرمنده (sharmande)}]{\fa{شرمنده} (sharmande) — ashamed; very sorry}

\label{lesson:FA-C482-sharmande}



\begin{quote}
Before the new one: say the Persian for hurt, then the Persian for a grudge.
\end{quote}

\subsection*{You'll want to know: \fa{شرمنده}}
\textbf{\fa{شرمنده}} (\emph{sharmande}) — \textquotedblleft{}ashamed; very sorry\textquotedblright{}.

\begin{quote}
\fa{او} \fa{از} \fa{رفتار} \fa{بدش} \fa{شرمنده} \fa{بود}.
\end{quote}

\begin{quote}
\emph{u az raftâr-e badash sharmande bud.}
\end{quote}

\begin{quote}
He was ashamed of his bad behaviour.
\end{quote}

Also a polite formula: \emph{sharmande-am} 'I'm so sorry (to trouble you)'. Built on \emph{sharm} 'shame'.

\begin{cousinweb}
\emph{sharmande} is built on the native Persian noun \emph{sharm}, 'shame', with the ending \emph{-ande}. EtymDB lists Hindi \emph{sharmindā} as a borrowing of it, but gives no source for the Persian word itself (EtymDB).
\end{cousinweb}

\begin{grammarlens}[title={What you've built}]
One more word for telling, explaining and describing at B1.
\end{grammarlens}

\subsection*{Your turn}
\begin{itemize}
\raggedright
  \item \emph{Say it:} \emph{sharmande}
  \item \emph{Say it:} \emph{sharmande}, once more
  \item \emph{Say it:} a sentence of your own with \emph{sharmande}
  \item \emph{Recall:} say the Persian for hurt, then the Persian for a grudge, then say \emph{sharmande} again
\end{itemize}

\begin{tcolorbox}[breakable,colback=teal!4,colframe=teal!35!black,title={Before you move on}]
What does \fa{شرمنده} mean? (\textquotedblleft{}Ashamed; very sorry\textquotedblright{}.) Say it once more.
\end{tcolorbox}
